\section*{References}

Abdelazeem, B., Hamdallah, A., Rizk, M. A., Abbas, K. S., El-Shahat, N. A., Manasrah, N., \& Mostafa, M. R. (2023). Does usage of monetary incentive impact the involvement in surveys? A systematic review and meta-analysis of 46 randomized controlled trials. \textit{PLOS ONE, 18}(1), e0279128. https://doi.org/10.1371/journal.pone.0279128

Aday, L. A. (1996). \textit{Designing and conducting health surveys: A comprehensive guide} (2nd ed.). Jossey-Bass.

Angur, M. G., \& Nataraajan, R. (1995). Do source of mailing and monetary incentives matter in international industrial mail surveys? \textit{Industrial Marketing Management, 24}(5), 351–357. https://doi.org/10.1016/0019-8501(95)00016-4

Babbie, E. R. (1990). \textit{Survey research methods} (2nd ed.). Cengage Learning.

Backstrom, C. H., \& Hursh-César, G. (1963). \textit{Survey research}. Northwestern University Press.

Balakrishnan, S., Chawla, S. K., Smith, M. F., \& Michalski, B. P. (1992). Mail survey response rates using a lottery prize giveaway incentive. \textit{Journal of Direct Marketing, 6}(3), 54–59. https://doi.org/10.1002/dir.4000060308

Baron, G., De Wals, P., \& Milord, F. (2001). Cost-effectiveness of a lottery for increasing physicians’ responses to a mail survey. \textit{Evaluation \& the Health Professions, 24}(1), 47–52. https://doi.org/10.1177/01632780122034777

Bee, A., \& Rothbaum, J. (2024). Using administrative data to evaluate nonresponse bias in the 2024 Current Population Survey Annual Social and Economic Supplement. \textit{U.S. Census Bureau}. https://www.census.gov/newsroom/blogs/research-matters/2024/09/administrative-data-nonresponse-bias-cps-asec.html

Blythe, B. J. (1986). Increasing mailed survey responses with a lottery. \textit{Social Work Research and Abstracts, 22}(3), 18–19. https://doi.org/10.1093/swra/22.3.18

Bosnjak, M., \& Tuten, T. L. (2003). Prepaid and promised incentives in web surveys: An experiment. \textit{Social Science Computer Review, 21}(2), 208–217. https://doi.org/10.1177/0894439303021002006

Brehm, J. (1994). Stubbing our toes for a foot in the door? Prior contact, incentives, and survey response. \textit{International Journal of Public Opinion Research, 6}(1), 45–63. https://doi.org/10.1093/ijpor/6.1.45

Brennan, M., Seymour, P., \& Gendall, P. (1993). The effectiveness of monetary incentives in mail surveys: Further data. \textit{Journal of the Market Research Society, 35}, 311–326.

Brick, J. M., \& Williams, D. (2013). Explaining rising nonresponse rates in cross-sectional surveys. \textit{The Annals of the American Academy of Political and Social Science, 645}(1), 36–59. https://doi.org/10.1177/0002716212456834

Church, A. H. (1993). Estimating the effect of incentives on mail survey response rates: A meta-analysis. \textit{Public Opinion Quarterly, 57}(1), 62–79. https://doi.org/10.1086/269355

Cobanoglu, C., \& Cobanoglu, N. (2003). The effect of incentives in web surveys: Application and ethical considerations. \textit{International Journal of Market Research, 45}(4), 475–488. https://doi.org/10.1177/147078530304500406

Conn, K., Mo, C. H., \& Purohit, B. (2025). Differential efficacy of survey incentives across contexts: Experimental evidence from Australia, India, and the United States. \textit{Political Science Research and Methods, 13}(4), 1055–1064. https://doi.org/10.1017/psrm.2024.53

Cook, C., Heath, F., \& Thompson, R. L. (2000). A meta-analysis of response rates in web- or internet-based surveys. \textit{Educational and Psychological Measurement, 60}(6), 821–836. https://doi.org/10.1177/00131640021970934

Coryn, C. L. S., Becho, L. W., Westine, C. D., Mateu, P. F., Abu-Obaid, R. N., Hobson, K. A., Schröter, D. C., Dodds, E. L., Vo, A. T., \& Ramlow, M. (2020). Material incentives and other factors associated with response rates to internet surveys. \textit{American Journal of Evaluation, 41}(2), 277–296. https://doi.org/10.1177/1098214018818371

Cottrell, E., Roddy, E., Rathod, T., Thomas, E., Porcheret, M., \& Foster, N. E. (2015). Maximising response from GPs to questionnaire surveys. \textit{BMC Medical Research Methodology, 15}, 3. https://doi.org/10.1186/1471-2288-15-3

Curtin, R., Presser, S., \& Singer, E. (2005). Changes in telephone survey nonresponse over the past quarter century. \textit{Public Opinion Quarterly, 69}(1), 87–98. https://doi.org/10.1093/poq/nfi002

Denton, J. J., \& Tsai, C.-Y. (1991). Two investigations into the influence of incentives and subject characteristics on mail survey responses in teacher education. \textit{Journal of Experimental Education, 59}(4), 352–366.

Deutskens, E., de Ruyter, K., Wetzels, M., \& Oosterveld, P. (2004). Response rate and response quality of internet-based surveys. \textit{Marketing Letters, 15}(1), 21–36. https://doi.org/10.1023/B:MARK.0000021968.86465.00

Drummond, F. J., O’Leary, E., O’Neill, C., Burns, R., \& Sharp, L. (2014). “Bird in the hand” cash was more effective than prize draws. \textit{Journal of Clinical Epidemiology, 67}(2), 228–231. https://doi.org/10.1016/j.jclinepi.2013.08.016

Edwards, P., Roberts, I., Clarke, M., DiGuiseppi, C., Pratap, S., \& Wentz, R. (2002). Increasing response rates to postal questionnaires: Systematic review. \textit{BMJ, 324}(7347), 1183. https://doi.org/10.1136/bmj.324.7347.1183

Finsen, V., \& Storeheier, A. H. (2006). Scratch lottery tickets are a poor incentive. \textit{BMC Medical Research Methodology, 6}, 19. https://doi.org/10.1186/1471-2288-6-19

Fox, R. J., Crask, M. R., \& Kim, J. (1988). Mail survey response rate: A meta-analysis. \textit{Public Opinion Quarterly, 52}(4), 467–491. https://doi.org/10.1086/269125

Göritz, A. S. (2006). Incentives in web studies: Methodological issues and a review. \textit{International Journal of Internet Science, 1}(1), 58–70.

Groves, R. M., Singer, E., \& Corning, A. (2000). Leverage-saliency theory of survey participation. \textit{Public Opinion Quarterly, 64}(3), 299–308. https://doi.org/10.1086/317990

Holbrook, A. L., Krosnick, J. A., \& Pfent, A. (2007). The causes and consequences of response rates in surveys. In \textit{Advances in telephone survey methodology} (pp. 499–528). Wiley.

Schwarz, N., \& Clore, G. L. (1996). Feelings and phenomenal experiences. In E. T. Higgins \& A. Kruglanski (Eds.), \textit{Social psychology: Handbook of basic principles} (pp. 433–465). Guilford Press.

Singer, E., \& Ye, C. (2013). The use and effects of incentives in surveys. \textit{The Annals of the American Academy of Political and Social Science, 645}(1), 112–141. https://doi.org/10.1177/0002716212458082

Tversky, A., \& Kahneman, D. (1981). The framing of decisions and the psychology of choice. \textit{Science, 211}(4481), 453–458. https://doi.org/10.1126/science.7455683

Yu, J., \& Cooper, H. (1983). A quantitative review of research design effects on response rates. \textit{Journal of Marketing Research, 20}(1), 36–44. https://doi.org/10.1177/002224378302000105


